% =============================================================================
%  MATERIAL DE ESTUDO - DESAFIO 08 (PDS 2026.2)
%  Discretização de Filtros IIR por Invariância ao Impulso e
%  Análise de Aliasing Espectral
%
%  Revisão dos conceitos necessários, de Sinais e Sistemas até PDS,
%  com 3 exercícios resolvidos por assunto (48 exercícios no total).
%
%  Compilar:
%      pdflatex material_estudo_desafio08.tex     (2x)
%
%  Pacotes exigidos (repositório TinyTeX desta máquina):
%      beamer pgf translator booktabs babel-portugues lm
%      amsmath amsfonts graphics tools geometry hyperref xcolor amscls
%
%  Todos os números deste documento foram calculados e conferidos em
%  material/gabarito_numerico.py  ->  material/gabarito_saida.txt
% =============================================================================
\documentclass[11pt,a4paper]{article}

\usepackage[utf8]{inputenc}
\usepackage[T1]{fontenc}
\usepackage[portuguese]{babel}
\usepackage{lmodern}
\usepackage{amsmath,amssymb,amsthm}
\usepackage{geometry}
\geometry{a4paper,top=2.4cm,bottom=2.4cm,left=2.4cm,right=2.4cm}
\usepackage{xcolor}
\usepackage{graphicx}
\usepackage{booktabs}
\usepackage{array}
\usepackage{hyperref}

\definecolor{azulPDS}{RGB}{31,78,121}
\definecolor{vermPDS}{RGB}{192,0,0}
\definecolor{verdePDS}{RGB}{46,139,87}
\definecolor{cinzaClaro}{RGB}{245,246,248}
\definecolor{amareloClaro}{RGB}{255,250,235}
\definecolor{azulClaro}{RGB}{238,244,250}

\hypersetup{colorlinks=true,linkcolor=azulPDS,urlcolor=azulPDS,
            citecolor=azulPDS,pdftitle={Material de Estudo - Desafio 08 - PDS 2026.2}}

% sumario apenas com as secoes (e subsecoes), sem os \paragraph
\setcounter{tocdepth}{1}
\setcounter{secnumdepth}{2}

\graphicspath{{../figuras/}}

% ------------------------------------------------------------------ ambientes
\newtheorem{teo}{Teorema}[section]
\newtheorem{prop}[teo]{Propriedade}
\newtheorem{exem}[teo]{Exemplo}
\theoremstyle{definition}
\newtheorem{defi}[teo]{Definição}
\theoremstyle{remark}
\newtheorem*{obs}{Observação}

\newcounter{ex}[section]
\newcommand{\exer}[1]{%
  \refstepcounter{ex}%
  \medskip\noindent\textbf{\color{azulPDS}Exercício \thesection.\theex\ --- #1}%
  \par\nopagebreak\smallskip}

\newcommand{\resp}[1]{%
  \begin{center}%
    \fcolorbox{vermPDS}{amareloClaro}{\parbox{0.94\textwidth}{\small
    \textbf{\color{vermPDS}Resposta:} #1}}%
  \end{center}}

\newcommand{\ideia}[1]{%
  \begin{center}%
    \fcolorbox{azulPDS}{azulClaro}{\parbox{0.94\textwidth}{\small
    \textbf{\color{azulPDS}Ideia central:} #1}}%
  \end{center}}

\newcommand{\jornal}[1]{\textcolor{cinzaClaro}{#1}}

\title{\vspace{-1.5cm}\textbf{\color{azulPDS}Material de Estudo --- Desafio 08}\\[2mm]
\large Discretização de Filtros IIR por Invariância ao Impulso e\\
Análise de Aliasing Espectral}
\author{PDS 2026.2 --- Departamento de Eletrônica e Sistemas (UFPE)\\
\small Revisão de Sinais e Sistemas $\rightarrow$ Processamento Digital de Sinais\\
\small \emph{48 exercícios resolvidos (3 por assunto)}}
\date{\today}

\begin{document}

\maketitle
\thispagestyle{empty}
\vspace{-1.5cm}

\begin{center}
\fcolorbox{azulPDS}{cinzaClaro}{\parbox{0.92\textwidth}{\small
\textbf{Como usar este material.} Cada capítulo tem (i) a teoria mínima
necessária, com as equações que você precisa saber \emph{de cor}, e (ii)
\textbf{três exercícios resolvidos} no mesmo estilo dos desafios da disciplina.
Todos os números foram obtidos por script (\texttt{gabarito\_numerico.py}), com a
saída em \texttt{gabarito\_saida.txt}: se um valor divergir do seu cálculo,
desconfie do seu cálculo. Os capítulos 4, 6, 7, 10, 12, 14, 15 e 16 são
diretamente o ``DNA'' do Desafio 08.}}
\end{center}

\tableofcontents
\newpage

\newcommand{\capitulo}[1]{%
  \IfFileExists{#1}{\input{#1}}{%
  \section*{[Capítulo ainda não escrito]}}}

\capitulo{capitulos/cap01_sinais_sistemas}
\capitulo{capitulos/cap02_slit_convolucao}
\capitulo{capitulos/cap03_laplace}
\capitulo{capitulos/cap04_segunda_ordem}
\capitulo{capitulos/cap05_fourier}
\capitulo{capitulos/cap06_amostragem}
\capitulo{capitulos/cap07_reconstrucao_poisson}
\capitulo{capitulos/cap08_discretos_lti}
\capitulo{capitulos/cap09_transformada_z}
\capitulo{capitulos/cap10_ft_digital}
\capitulo{capitulos/cap11_dft_janelas}
\capitulo{capitulos/cap12_fase_grupo}
\capitulo{capitulos/cap13_fir_fase_linear}
\capitulo{capitulos/cap14_estruturas_quantizacao}
\capitulo{capitulos/cap15_invariancia_impulso}
\capitulo{capitulos/cap16_bilinear}

\appendix
\capitulo{capitulos/apA_formulario}
\capitulo{capitulos/apB_desafio08_resolvido}
\capitulo{capitulos/apC_verificacao}

\end{document}
